\documentclass[11pt,a4paper]{article}
\usepackage[T1]{fontenc}
\usepackage{lmodern,microtype}
\usepackage[margin=26mm,headheight=15pt]{geometry}
\usepackage{amsmath,amssymb,amsthm,mathtools}
\usepackage{booktabs,array,longtable,enumitem}
\usepackage[dvipsnames]{xcolor}
\usepackage{fancyhdr}
\usepackage{hyperref}
\hypersetup{colorlinks=true,linkcolor=MidnightBlue,citecolor=MidnightBlue,urlcolor=MidnightBlue,pdftitle={Sharp Subexponential Reversion},pdfauthor={ChatGPT, prepared for Vladimir Reshetnikov}}
\usepackage[nameinlink,noabbrev]{cleveref}
\setlength{\parskip}{4pt}
\setlength{\parindent}{0pt}
\setlength{\emergencystretch}{2em}
\pagestyle{fancy}\fancyhf{}
\fancyhead[L]{\small Sharp subexponential reversion}
\fancyhead[R]{\small Research article / 29 September 2026}
\fancyfoot[C]{\thepage}
\newtheorem{theorem}{Theorem}[section]
\newtheorem{proposition}[theorem]{Proposition}
\newtheorem{lemma}[theorem]{Lemma}
\newtheorem{corollary}[theorem]{Corollary}
\theoremstyle{definition}
\newtheorem{definition}[theorem]{Definition}
\newtheorem{example}[theorem]{Example}
\theoremstyle{remark}
\newtheorem{remark}[theorem]{Remark}
% ed. (2026-09-29): unnumbered environment for ProveIt editorial notes; remarks share the theorem counter, so the notes must not be numbered.
\newtheorem*{ednote}{Editorial note (ProveIt, 2026-09-29)}
\newcommand{\C}{\mathbb C}
\newcommand{\N}{\mathbb N}
\newcommand{\R}{\mathbb R}
\newcommand{\ee}{\mathrm e}
\newcommand{\typ}{\tau_s}
\newcommand{\coef}[2]{[z^{#1}]\,#2}
\newcommand{\inv}{^{\langle-1\rangle}}
\newcommand{\Prob}{\mathbb P}
\newcommand{\Exp}{\mathbb E}
\newcommand{\abs}[1]{\lvert#1\rvert}
\newcommand{\set}[1]{\{#1\}}
\newcommand{\commit}{\texttt{5c695cfd\,f70a8b6c\,91bd5b2c\,1e99e3ba\,ecd5a864}}
\DeclareMathOperator{\Pois}{Pois}
\DeclareMathOperator{\Var}{Var}
\DeclareMathOperator{\supp}{supp}
\numberwithin{equation}{section}
\begin{document}
% ed. (2026-09-29): no PDF page anchor on the title page (the delivered build had a duplicate destination page.1)
\hypersetup{pageanchor=false}
\begin{titlepage}
\vspace*{10mm}
{\color{MidnightBlue}\Large ProveIt research continuation\par}
\vspace{12mm}
{\Huge\bfseries Sharp Subexponential\par Reversion\par}
\vspace{7mm}
{\Large Exact extremal cost, a Gevrey transition,\par and composition of two divergent series\par}
\vspace{12mm}
{\large Prepared for Vladimir Reshetnikov\par}
{\large ChatGPT\par}
{29 September 2026\par}
\vspace{12mm}
\begin{abstract}
Exact preservation of a weighted coefficient type under formal reversion
is already established in the inspected ProveIt corpus. It does not
quantify the subexponential amplification at the critical coefficient
radius. We address that subsequent question by solving a coefficientwise
extremal problem for two factorial normalizations. A single negative
series attains the worst inverse coefficient simultaneously at every
degree. For its normalized inverse amplification $R_n$, we prove
$\log R_n\sim Cn^{1-s}$ for every $0<s<1$, a multiplicative equivalent
$R_n\sim\exp(Cn^{1-s})$ for $s>1/2$, and the exact additional constant
at $s=1/2$. A uniform equivalent yields the crossover
$R_n\to\exp(C\ee^{-t})$ when $s=1+t/\log n$. Consequently, fixed-radius
factorial coefficient balls are stable under inversion exactly for
$s\ge1$, despite exact type preservation at every positive order.

A second result answers the Gevrey version of the repository's question
about two divergent composition arguments: their types obey a sharp
maximum inequality, with equality whenever the two competing types are
unequal. Finally, an explicit positive-coefficient series has a full
normalized root limit although its inverse has infinitely many zero
coefficients. All proofs are supplied. Finite exact checks and numerical
diagnostics are distinguished from proofs and from Lean verification.
\end{abstract}
\vfill
\small
\textbf{Scope and status.} Conventional mathematical proofs and reproducible
finite calculations. No Lean formalization, analytic summation theorem,
or worldwide priority claim is asserted. The contribution is measured
against explicitly identified questions in the inspected repository.
\end{titlepage}
% ed. (2026-09-29): page anchors resume after the title page
\hypersetup{pageanchor=true}
\tableofcontents
\clearpage

\section{Research target and contribution boundaries}\label{sec:target}

\subsection{Which repository questions are actually still being extended?}
The source for this continuation is the ProveIt package
\emph{Sharp Weighted Type Under Formal Reversion: zero-loss inversion,
exact Borel radii, and non-D-finiteness in countable exponential feedback}
\cite{repo,weighted}. In its section on further research, the subsection
``Optimal subexponential cost of reversion'' asks for the smallest
universal overhead at fixed weight and fixed input constants, naming
$M_n=(n!)^s$, $0<s<1$, as a first target. The subsection ``Two divergent
composition arguments'' asks for a filtered composition estimate and a
noncancellation condition when both arguments are divergent. These are
the two concrete questions treated here.

The inspection used the main-branch snapshot identified by
\begin{center}\commit\end{center}
The relevant source file is
\begin{quote}\small\ttfamily
Analysis/Transseries/docs/series-and-transseries/\\
Sharp\_Weighted\_Type\_Formal\_Reversion/article.tex.
\end{quote}
Its research questions occur in the inspected range 1418--1515; the
individual subsection headings, rather than mutable line numbers, are
the durable locators. The directory inventory and editorial notes were
also consulted \cite{inventory}.

There is a significant duplication trap. The amplitude--slope package
\cite{amplitude} asks whether inverse and forward Gevrey types agree.
The current inventory explicitly observes that the weighted-type theorem
already answers that question. Accordingly, merely reproving equality of
these types would not be a new repository contribution. The present work
moves below the $n$th-root scale and treats two genuinely divergent
composition inputs. The older type-invariance statement is reproved only
as a short, self-contained tool.

\subsection{Results and exact extent of the resolution}
The principal result, \cref{thm:main}, identifies the exact finite-degree
extremum and its leading subexponential asymptotics. For $0<s<1$, the
coefficient in the exponent is determined, not just the power of $n$.
For $s\ge1/2$, the statements give a full multiplicative equivalent,
including the separate constant at the boundary $s=1/2$. For $0<s<1/2$,
a complete multiplicative equivalent remains open here: further
subexponential terms cannot in general be discarded.

The maximum law in \cref{thm:composition} settles the factorial-weight
version of the second question. It does not settle the same problem for
every superexponential log-convex weight. Equal-type cancellation is
unavoidable; unequal-type inputs supply an explicit sufficient
noncancellation condition.
% ed. (2026-09-29): the general-weight version is settled by a companion package of the same batch
\begin{ednote}
The companion package \path{Exact_Weighted_Type_Beyond_Log_Convexity/}
(filed in the same batch, written independently from an earlier snapshot)
proves the same maximum law, with equality for unequal types, for every
positive weight that is subexponentially equivalent to a supermultiplicative
weight with roots tending to infinity, in particular for every log-convex
weight with $M_n^{1/n}\to\infty$ (its Corollary~\texttt{cor:scalarscaling}
and Theorem~\texttt{thm:composition}). It also shows that at equal types $T$
every composite type in $[0,T]$ occurs (its
Theorem~\texttt{thm:cancellation}). The results of
Section~\ref{sec:composition} are the factorial case of that law, proved
here by a different majorant: one theorem with two independent proofs. The
fixed-ball results of Sections~\ref{sec:main}--\ref{sec:type} are not in
that package.
\end{ednote}

The positive-forward counterexample in \cref{thm:rootcounter} explains
why exact type alone cannot be strengthened to a full inverse root limit,
even after assuming positivity and a full root limit on the forward
side. It is a general obstruction, not a claim about the signs of a
particular exponential-feedback inverse.

\subsection{Relation to established literature}
Formal Lagrange inversion is classical; Gessel's survey supplies a
systematic account of its coefficient forms and proofs \cite{gessel}.
Composition and inverse stability in ultradifferentiable classes have
been studied by Rainer and Schindl \cite{rscomposition,rsstability}.
Borinsky develops composition and inversion in a factorially divergent
coefficient-asymptotic algebra \cite{borinsky}. These precedents are
relevant: neither Gevrey closure nor the usefulness of a single large
coefficient should be advertised as newly discovered here.

Our object is a sharp extremum over a \emph{fixed formal coefficient
ball}, including its dependence on the order $s$, rather than stability
of a function class after increasing an unspecified exponential constant.
The proof is elementary and includes the exact constants needed for this
object. The literature check was targeted, not an exhaustive priority
search; the results should be independently reviewed before publication
or incorporation as established repository theorems.

\begin{table}[htbp]
\centering\small
\begin{tabular}{p{.34\textwidth}p{.59\textwidth}}
\toprule
Assertion & Status in this article\\
\midrule
Coefficientwise extremal inverse & Exact finite identity and proof.\\
Optimal logarithmic cost, $0<s<1$ & Upper and lower bounds with the same constant.\\
Multiplicative cost, $s>1/2$ & Proved, locally uniformly in $s$ and $C$.\\
Boundary $s=1/2$ & Proved with its additional constant.\\
Two divergent composition inputs & Sharp type inequality; equality at unequal types.\\
General type invariance & Existing repository result; factorial case reproved.\\
Full feedback inverse asymptotics & Not inferred from a universal upper envelope.\\
Computations & Exact finite tests and non-certified floating-point diagnostics.\\
\bottomrule
\end{tabular}
\caption{Dependency and claim ledger. None of the mathematical results in this article is presented as Lean-verified.}
\end{table}

\section{The extremal problem and the main theorem}\label{sec:main}

\subsection{Two normalizations, kept explicit}
Fix $s>0$, $A>0$, $C>0$, and $\nu\in\{0,1\}$. For $j\ge1$, define
\begin{equation}
 w_j=w_j(s,\nu)=
 \left(\frac{(j+\nu)!}{(1+\nu)!}\right)^s.
 \label{eq:weights}
\end{equation}
In particular,
\begin{equation}
 w_1=1,\qquad w_2=(\nu+2)^s,\qquad
 \frac{w_{j+1}}{w_j}=(j+\nu+1)^s.
 \label{eq:ratios}
\end{equation}
The symbol $w_0$ is not needed. When it is useful in a computational
recurrence, we set $w_0=1$ as a separate convention.

\begin{definition}[A critical-radius coefficient ball]
Let $\mathcal B_{s,A,C}^{(\nu)}$ be the set of all formal series
\begin{equation}
 f(z)=z+\sum_{j\ge1}a_{j+1}z^{j+1}\in\C[[z]],\qquad
 |a_{j+1}|\le C A^j w_j.
 \label{eq:ball}
\end{equation}
For $g=f\inv$, define the normalized extremal amplification
\begin{equation}
 R_n(s,C,\nu)=
 \sup_{f\in\mathcal B_{s,A,C}^{(\nu)}}
 \frac{|[z^n]f\inv|}{C A^{n-1}w_{n-1}},\qquad n\ge2.
 \label{eq:extremum}
\end{equation}
The theorem below shows that this expression is independent of $A$.
The inverse is compositional, not the reciprocal $1/f$.
\end{definition}

The case $\nu=0$ uses $(n-1)!^s$ for the $n$th coefficient. The case
$\nu=1$ uses $(n!/2)^s$. Thus the conventional ball
\begin{equation}
 |a_n|\le B A^{n-1}(n!)^s\quad(n\ge2)
 \label{eq:standardball}
\end{equation}
is exactly \eqref{eq:ball} with $\nu=1$ and $C=2^sB$.
The two weights have identical $n$th-root types but different fixed-ball
constants. This difference matters in every sharp subexponential formula.

\begin{theorem}[Sharp extremal reversion cost]\label{thm:main}
Let $s,A,C,\nu$ be as above. Define
\begin{equation}
 F_*(z)=z-C\sum_{j\ge1}A^j w_j z^{j+1},\qquad G_*=F_*\inv.
 \label{eq:extremizer}
\end{equation}
All coefficients of $G_*$ after the constant term are strictly positive.
The same series $F_*$ attains the supremum in \eqref{eq:extremum} for
every $n$. Moreover:
\begin{enumerate}[label=(\roman*),leftmargin=9mm]
\item For every fixed $0<s<1$,
\begin{equation}
 \lim_{n\to\infty}\frac{\log R_n(s,C,\nu)}{n^{1-s}}=C.
 \label{eq:logmain}
\end{equation}
\item For every fixed $s>1/2$,
\begin{equation}
 R_n(s,C,\nu)=\exp(Cn^{1-s})\,(1+o(1)).
 \label{eq:multmain}
\end{equation}
The $o(1)$ is uniform for $s$ in a compact subset of $(1/2,\infty)$
and $C$ in a compact subset of $(0,\infty)$.
\item At $s=1/2$,
\begin{equation}
 R_n(1/2,C,\nu)=
 \exp\left(C\sqrt n+\frac{3C^2}{4}+C\sqrt{\nu+2}\right)(1+o(1)).
 \label{eq:halfmain}
\end{equation}
\item Uniformly for $t$ in every bounded real interval,
\begin{equation}
 R_n\left(1+\frac{t}{\log n},C,\nu\right)
 \longrightarrow \exp(C\ee^{-t}).
 \label{eq:crossover}
\end{equation}
\end{enumerate}
\end{theorem}

The theorem gives the following sharp classification:
\begin{equation}
 \begin{array}{c|c}
 \text{order} & \text{extremal amplification}\\ \hline
 s>1 & R_n\longrightarrow1\\
 s=1 & R_n\longrightarrow\ee^C\\
 0<s<1 & R_n\longrightarrow\infty,\quad
          \log R_n\sim Cn^{1-s}.
 \end{array}
 \label{eq:phase}
\end{equation}
At $0<s<1$, this is an amplification at the \emph{same} $A$; it is
still subexponential and hence invisible to the coefficient type.
% ed. (2026-09-29): relation to the companion package's envelope asymptotic
\begin{ednote}
With $A=1$, $G_*$ is the positive extremal inverse $H_C$ of the companion
package \path{Exact_Weighted_Type_Beyond_Log_Convexity/} for the
excess-degree weight $N_j=w_j$, which is log-convex with $w_0=1$ and hence
its own supermultiplicative envelope. That package proves $\log R_n=o(n)$
for every weight with superexponential roots (its
Theorem~\texttt{thm:envelopeintro}) and asks for the optimal overhead;
Theorem~\ref{thm:main} answers that question for factorial weights.
\end{ednote}

\begin{corollary}[The unshifted factorial convention]\label{cor:unshifted}
For the ball \eqref{eq:standardball}, normalize inverse coefficients by
$B A^{n-1}(n!)^s$ and call their extremum $\widetilde R_n$.
Then, for $0<s<1$,
\[
 \log\widetilde R_n\sim 2^s B n^{1-s}.
\]
For $s>1/2$, $\widetilde R_n\sim\exp(2^sBn^{1-s})$. At $s=1/2$,
\[
 \widetilde R_n\sim
 \exp\left(\sqrt2 B\sqrt n+\frac32 B^2+\sqrt6 B\right).
\]
In particular, $\widetilde R_n\to\ee^{2B}$ when $s=1$.
\end{corollary}
\begin{proof}
Use $\nu=1$ and $C=2^sB$ in \cref{thm:main}.
\end{proof}

\section{Finite coefficient algebra and uniform convolution bounds}
\label{sec:finite}

\subsection{An exact simultaneous extremizer}
For $f=z(1+v)$ with $v\in z\C[[z]]$, formal Lagrange inversion gives
\begin{equation}
 [z^n]f\inv=\frac1n[z^{n-1}](1+v(z))^{-n}.
 \label{eq:lagrange}
\end{equation}
Every coefficient is a finite sum; no radius of convergence is needed.
Since
\[
 (1+v)^{-n}=\sum_{k\ge0}(-1)^k\binom{n+k-1}{k}v^k,
\]
taking absolute values and then replacing each $|[z^j]v|$ by its upper
bound in \eqref{eq:ball} gives equality when all coefficients of $v$ are
negative and saturate the bounds.

For $m\ge k\ge1$, put
\begin{equation}
 P(m,k)=\sum_{\substack{j_1+\cdots+j_k=m\\j_i\ge1}}
                   \prod_{i=1}^k w_{j_i}.
 \label{eq:convolution}
\end{equation}
Also set $P(0,0)=1$ and $P(m,0)=0$ for $m>0$.
Then the exact value of the extremum is
\begin{equation}
 R_n=\frac{1}{n w_{n-1}}
 \sum_{k=1}^{n-1}\binom{n+k-1}{k}C^{k-1}P(n-1,k).
 \label{eq:exactR}
\end{equation}
All powers of $A$ cancel because the indices in each composition sum to
$n-1$. The $k=1$ summand is exactly $1$, so $R_n\ge1$. This proves
all finite extremizer assertions of \cref{thm:main}.

\subsection{A factorial convolution estimate}
\begin{lemma}[Uniform two-part convolution]\label{lem:twoconv}
There is a finite $K=K(s,\nu)\ge1$ such that
\begin{equation}
 \sum_{j=1}^{m-1} w_jw_{m-j}\le K w_{m-1}\qquad(m\ge2).
 \label{eq:twoconv}
\end{equation}
The constant can be chosen bounded when $s$ varies in any compact subset
of $(0,\infty)$. Consequently,
\begin{equation}
 P(m,k)\le K^{k-1}w_{m-k+1}\qquad(m\ge k\ge1).
 \label{eq:manyconv}
\end{equation}
\end{lemma}
\begin{proof}
For $1\le j\le m/2$,
\[
 \binom mj=\frac mj\binom{m-1}{j-1}\ge\frac mj\,2^{j-1}.
\]
Indeed, each factor in the product for $\binom{m-1}{j-1}$ is at least
$2$ when $m\ge2j$. For $\nu=0$ this gives
\[
 \frac{w_jw_{m-j}}{w_{m-1}}
 =\left(\frac m{\binom mj}\right)^s
 \le\left(\frac j{2^{j-1}}\right)^s.
\]
For $\nu=1$ the ratio is
\[
 \left(\frac{(j+1)(m-j+1)}{2\binom mj}\right)^s
 \le\left(\frac{j(j+1)}{2^j}\right)^s.
\]
Reflect about $m/2$ and sum. Twice the corresponding infinite sum is a
valid $K$ after increasing it to at least $1$. Both sums converge
uniformly on compact positive $s$-intervals.

To prove \eqref{eq:manyconv}, use induction on $k$. Splitting off the
first part gives
\[
 P(m,k)\le K^{k-2}\sum_{j=1}^{m-k+1}w_jw_{m-k+2-j}
 \le K^{k-1}w_{m-k+1}.
\]
The case $k=1$ is an equality.
\end{proof}

Let $T_{n,k}$ denote the $k$th summand in \eqref{eq:exactR}. For
$k\le n/2$, \cref{lem:twoconv} implies
\begin{equation}
 T_{n,k}\le
 2\frac{(D n^{1-s})^{k-1}}{k!},\qquad
 D=2^{s+1}CK.
 \label{eq:poissontailbound}
\end{equation}
In fact, $\binom{n+k-1}{k}\le(2n)^k/k!$, and the ratio
$w_{n-k}/w_{n-1}$ contains $k-1$ factors bounded by $(n/2)^{-s}$.
For $k>n/2$, at least $\lfloor n/2\rfloor-1$ such factors remain,
whereas the binomial coefficient is at most $4^n$. Thus
\begin{equation}
 \sum_{k>n/2}T_{n,k}
 \le\exp\left(-\frac{s}{2}n\log n+O(n)\right).
 \label{eq:longtail}
\end{equation}
Here and below constants are allowed to depend on fixed $s,C,\nu$.
They are uniform on compact positive $s,C$ sets when that is asserted.
Together these estimates give
\begin{equation}
 1\le R_n\le 2\exp(D n^{1-s})+
            \exp\left(-\frac{s}{2}n\log n+O(n)\right).
 \label{eq:coarsecost}
\end{equation}
This already proves subexponential overhead. Determining its sharp
constant requires more than replacing every convolution by $K$.

\section{Localization at one large coefficient}\label{sec:giant}

A positive composition of $m=n-1$ is called \emph{localized at scale $h$}
if one of its parts is at least $m-h$, where $h<m/2$. Such a part is
unique. This terminology concerns the finite Lagrange sum only.

\begin{lemma}[Negligible configurations without a large part]
\label{lem:localization}
Let the contribution to \eqref{eq:exactR} of nonlocalized compositions
be $E_n(h)$. Then:
\begin{enumerate}[label=(\roman*),leftmargin=9mm]
\item For fixed $0<\epsilon<1/2$ and $h=\lfloor\epsilon m\rfloor$,
$E_n(h)\le\exp(-c_\epsilon n)$ for all sufficiently large $n$.
\item If
\begin{equation}
 \max(1-s,0)<\kappa<\alpha<1,\qquad h=\lfloor n^\alpha\rfloor,
 \label{eq:localpowers}
\end{equation}
then $E_n(h)=o(1)$. More precisely it is bounded by the sum of
$\exp(-c n^\alpha\log n)$ and $\exp(-c' n^\kappa\log n)$.
\end{enumerate}
\end{lemma}
\begin{proof}
First restrict to compositions of length at most $t=o(n)$. If every
part is at most $m-h$ and $m-h\ge m/2$, elementary factorial
concentration gives
\begin{equation}
 \prod_i j_i!\le(m-h)!h!.
 \label{eq:capconcentration}
\end{equation}
To see this, allow zero parts and transfer units from a smaller positive
part to a larger one, without exceeding $m-h$. Each transfer weakly
increases the factorial product. A maximizer has at most two positive
parts, of sizes $m-h$ and $h$.

For $\nu=0$, the normalized weight product is therefore at most
$\binom mh^{-s}$. For $\nu=1$, there is an additional factor bounded
by
\[
 \left(\prod_i(j_i+1)\right)^s
 \le (1+m/k)^{sk}
\]
for a composition of length $k$, by the arithmetic--geometric mean
inequality. The factors $2^{-sk}$ and the extra normalization in $w_m$
only improve this upper bound.

There are $\binom{m-1}{k-1}$ compositions of length $k$. Both this
count and $\binom{n+k-1}{k}$ have logarithm
$O(k\log(\ee n/k))$. The powers of $C$ have logarithm $O(k)$.
Summing over $k\le t$ consequently bounds this portion by
\begin{equation}
 \exp\left\{-s\log\binom mh+
       O\left(t\log\frac{\ee n}{t}+\log n\right)\right\}.
 \label{eq:nogiantentropy}
\end{equation}
The function $k\log(\ee n/k)$ is increasing for the relevant $k$.

Lengths $k>t$ are controlled by \eqref{eq:poissontailbound} and
\eqref{eq:longtail}. If $t/n^{1-s}\to\infty$ and $t\le n/2$, the
small-length tail is bounded, after harmless polynomial factors, by
\begin{equation}
 \exp\left\{-t\log\frac{t}{\ee D n^{1-s}}\right\}.
 \label{eq:tailrate}
\end{equation}
This follows either from successive-term ratios or from $k!\ge(k/\ee)^k$.

For (i), choose $t=\lfloor n/\sqrt{\log n}\rfloor$.
The error term in \eqref{eq:nogiantentropy} is $o(n)$, while
$\log\binom mh=nH(\epsilon)+o(n)$, where
$H(x)=-x\log x-(1-x)\log(1-x)>0$. The exponent in
\eqref{eq:tailrate} is at most $-(s+o(1))n\sqrt{\log n}$.
For (ii), take $t=\lfloor n^\kappa\rfloor$.
Then $\log\binom mh=(1-\alpha+o(1))n^\alpha\log n$ dominates
$O(n^\kappa\log n)$, and \eqref{eq:tailrate} has a negative constant
multiple of $n^\kappa\log n$. This proves both statements.
\end{proof}

\subsection{The exact localized sum}
Write $m=n-1$ and let $R_n^{\mathrm{loc}}(h)$ be the localized
contribution. Let the large part have size $m-d$, let $\ell$ be the
number of other parts, and let $m_j$ count the other parts of size $j$.
Then
\begin{equation}
 \ell=\sum_{j\ge1}m_j,\qquad d=\sum_{j\ge1}jm_j.
 \label{eq:counts}
\end{equation}
Choosing the position of the large part changes the Lagrange factor by
\[
 \frac{\ell+1}{n}\binom{n+\ell}{\ell+1}=\binom{n+\ell}{\ell}.
\]
Thus, exactly,
\begin{equation}
 R_n^{\mathrm{loc}}(h)=
 \sum_{\substack{m_1,m_2,\ldots\ge0\\d\le h}}
 \binom{n+\ell}{\ell}\frac{\ell!}{\prod_jm_j!}\,
 C^\ell\,\frac{w_{m-d}}{w_m}\prod_jw_j^{m_j}.
 \label{eq:localizedexact}
\end{equation}
Only $j\le h$ can occur. The all-zero multiplicity vector contributes
$1$. Formula \eqref{eq:localizedexact} is the source of all subsequent
constants.

% ed. (2026-09-29): \texorpdfstring added (the delivered build warned about math in the PDF bookmark)
\section{Optimal logarithmic amplification for every \texorpdfstring{$0<s<1$}{0 < s < 1}}
\label{sec:logcost}

\subsection{Lower bound from quadratic decorations}
In \eqref{eq:localizedexact}, retain only parts of size $1$ beside the
large part. Their contribution for a given $\ell$ is
\begin{equation}
 L_{n,\ell}=\binom{n+\ell}{\ell} C^\ell
             \frac{w_{n-1-\ell}}{w_{n-1}}.
 \label{eq:ones}
\end{equation}
Take $\ell=\lfloor Cn^{1-s}\rfloor$. It tends to infinity and is
$o(n)$. Uniform elementary product estimates give
\begin{align}
 \log\binom{n+\ell}{\ell}
    &=\ell\log n-\log(\ell!)+O(\ell^2/n+\ell/n),\\
 \log\frac{w_{n-1-\ell}}{w_{n-1}}
    &=-s\ell\log n+O(\ell^2/n+\ell/n).
 \label{eq:onesest}
\end{align}
Stirling's estimate then yields
\[
 \log L_{n,\ell}
 =\ell\log\frac{Cn^{1-s}}{\ell}+\ell+o(\ell)
 =Cn^{1-s}+o(n^{1-s}).
\]
For large $n$ the part $n-1-\ell$ is genuinely the unique large part,
so these are actual distinct summands, not a formal overcount. Hence
\begin{equation}
 \liminf_{n\to\infty}\frac{\log R_n}{n^{1-s}}\ge C.
 \label{eq:lowercost}
\end{equation}

\subsection{Upper bound with the same constant}
Fix $0<\epsilon<1/2$, let $h=\lfloor\epsilon m\rfloor$, and put
\begin{equation}
 t_n=(m+\nu-h+1)^{-s}.
\end{equation}
For $d\le h$, each of the $d$ factors in the factorial quotient gives
\[
 \frac{w_{m-d}}{w_m}\le t_n^d.
\]
Define the finite polynomial $H_h(t)=\sum_{j=1}^h w_jt^j$.
Extending the positive sum over the small parts and then over $\ell$
in \eqref{eq:localizedexact} gives
\begin{equation}
 R_n^{\mathrm{loc}}(h)
 \le\sum_{\ell\ge0}\binom{n+\ell}{\ell}
                  (C H_h(t_n))^\ell
 =\bigl(1-C H_h(t_n)\bigr)^{-n-1}.
 \label{eq:negativebinomial}
\end{equation}
The last identity is used only when $n$ is sufficiently large that
$C H_h(t_n)<1$.

For $1\le j<h$, successive terms of $H_h(t_n)$ have ratio
$(j+\nu+1)^s t_n$. Its maximum tends to
$(\epsilon/(1-\epsilon))^s<1$. Therefore
\[
 H_h(t_n)=t_n+O(t_n^2),\qquad
 t_n=(1-\epsilon)^{-s}n^{-s}(1+o(1)).
\]
Taking logarithms in \eqref{eq:negativebinomial} proves
\[
 \log R_n^{\mathrm{loc}}(h)
 \le C(1-\epsilon)^{-s}n^{1-s}+o(n^{1-s}).
\]
The nonlocalized contribution is exponentially small by
\cref{lem:localization}(i), whereas the localized contribution is at
least $1$. Consequently,
\[
 \limsup_{n\to\infty}\frac{\log R_n}{n^{1-s}}
 \le C(1-\epsilon)^{-s}.
\]
Letting $\epsilon\downarrow0$ and using \eqref{eq:lowercost} proves
\eqref{eq:logmain}. In particular, no universal bound
$\exp(o(n^{1-s}))$ is possible for this fixed ball, and no smaller
leading constant than $C$ can hold.

\section{Multiplicative asymptotics and the boundary constant}
\label{sec:multcost}

\subsection{A finite Poisson representation}
The probability notation used here is just a convenient normalization of
a finite positive sum. Fix a cutoff $h=o(n)$ and define
\begin{equation}
 \lambda_j=Cw_j n^{1-sj}\quad(1\le j\le h),\qquad
 \Lambda_n=\sum_{j=1}^h\lambda_j.
 \label{eq:intensities}
\end{equation}
Let $N_j$ be independent Poisson variables of means $\lambda_j$ and put
\begin{equation}
 L=\sum_{j=1}^h N_j,\qquad D_0=\sum_{j=1}^h jN_j.
 \label{eq:poissonstats}
\end{equation}
The subscript on $D_0$ distinguishes this random total size from the
constant $D$ in \eqref{eq:poissontailbound}.

For a multiplicity vector of length $\ell$ and size $d$, extract
$n^\ell$ from $(n+1)\cdots(n+\ell)$ and $n^{-sd}$ from the large
weight quotient in \eqref{eq:localizedexact}. If $d\le h$,
\begin{equation}
 \log\left\{\frac{(n+1)\cdots(n+\ell)}{n^\ell}
       \frac{w_{n-1-d}}{w_{n-1}}n^{sd}\right\}
 =O\left(\frac{h^2+h}{n}\right).
 \label{eq:uniformfactor}
\end{equation}
The constants are uniform on compact $s$-intervals.
Thus, when $h^2/n\to0$,
\begin{equation}
 R_n^{\mathrm{loc}}(h)
 =(1+o(1))\ee^{\Lambda_n}\Prob(D_0\le h).
 \label{eq:poissonfirst}
\end{equation}
No probabilistic limit theorem is hidden in this identity: it follows by
multiplying and dividing each summand by $\prod_j\ee^{-\lambda_j}$.

\subsection{Proof of the equivalent above one half}
Assume first that $s>1/2$. Choose
\[
 \max(1-s,0)<\kappa<\alpha<1/2,\qquad h=\lfloor n^\alpha\rfloor.
\]
Nonlocalized configurations contribute $o(1)$ by
\cref{lem:localization}(ii). For the intensities in
\eqref{eq:intensities}, successive terms have ratio
\[
 \frac{\lambda_{j+1}}{\lambda_j}
 =\left(\frac{j+\nu+1}{n}\right)^s.
\]
Uniformly for $j<h$, this tends to zero. Since $w_1=1$,
\begin{equation}
 \Lambda_n=Cn^{1-s}+O(n^{1-2s}),\qquad
 \Exp D_0=Cn^{1-s}+O(n^{1-2s}).
 \label{eq:poissonmeans}
\end{equation}
The error in the first formula tends to zero. The mean in the second is
$o(h)$, so Markov's inequality gives $\Prob(D_0\le h)\to1$.
Formula \eqref{eq:poissonfirst} proves \eqref{eq:multmain}.

For a compact interval of $s$-values with lower endpoint $s_0>1/2$,
choose $\kappa,\alpha$ once with
$\max(1-s_0,0)<\kappa<\alpha<1/2$. Every estimate just used is then
uniform, as are the convolution constants. Compact variation in $C>0$
is harmless. This proves the asserted local uniformity.

At $s_n=1+t/\log n$, $n^{1-s_n}=\ee^{-t}$ exactly. For bounded $t$,
$s_n$ eventually belongs to a fixed compact subinterval of
$(1/2,\infty)$. The uniform formula gives \eqref{eq:crossover}.

\subsection{Why the order-one-half constant is different}
At $s=1/2$, the typical number of small parts is of order $\sqrt n$.
Their squared total size divided by $n$ is no longer negligible. In
addition, parts of size $2$ have a nonvanishing Poisson mean.

Choose
\begin{equation}
 \frac12<\kappa<\alpha<\frac23,\qquad h=\lfloor n^\alpha\rfloor.
 \label{eq:halfcutoff}
\end{equation}
Localization again leaves an absolute $o(1)$ remainder. Expansion of the
two finite products in \eqref{eq:uniformfactor} now gives, uniformly for
$\ell\le d\le h$,
\begin{align}
 \log\frac{(n+1)\cdots(n+\ell)}{n^\ell}
   &=\frac{\ell(\ell+1)}{2n}+O(h^3/n^2),\\
 \log\left(\frac{w_{n-1-d}}{w_{n-1}}n^{sd}\right)
   &=\frac{s\{d^2+(1-2\nu)d\}}{2n}+O(h^3/n^2).
 \label{eq:secondproduct}
\end{align}
The linear terms are $o(1)$, and $h^3/n^2=o(1)$. Hence
\begin{equation}
 R_n^{\mathrm{loc}}(h)
 =\ee^{\Lambda_n}\Exp\left[
   \exp\left(\frac{L^2+sD_0^2}{2n}+o(1)\right)
   \boldsymbol1_{\{D_0\le h\}}\right],
 \label{eq:secondpoisson}
\end{equation}
where the $o(1)$ inside the expectation is uniform on the indicated event.

Here
\begin{equation}
 \lambda_1=C\sqrt n,\qquad \lambda_2=C\sqrt{\nu+2},\qquad
 \sum_{j=3}^h j^r\lambda_j=O(n^{-1/2})\quad(r=0,1,2).
 \label{eq:halfmeans}
\end{equation}
It follows that
\begin{equation}
 \Lambda_n=C\sqrt n+C\sqrt{\nu+2}+o(1),\qquad
 \frac{L}{\sqrt n},\frac{D_0}{\sqrt n}\longrightarrow C
 \quad\text{in probability}.
 \label{eq:halfprob}
\end{equation}
For example, these convergences follow directly from the corresponding
means and variances, which are $C\sqrt n+O(1)$.

We must justify passage to the expectation of an exponential; convergence
in probability alone would not suffice. Fix $\theta>0$. Successive
intensities after multiplication by $\ee^{\theta j}$ still have ratio
$o(1)$, so
\begin{equation}
 \Exp\ee^{\theta D_0}
 =\exp\left(C(\ee^\theta-1)\sqrt n+O(1)\right).
 \label{eq:chernoff}
\end{equation}
On $D_0\le h$, we have $L\le D_0$ and
\[
 \frac{L^2+sD_0^2}{2n}\le\frac{(1+s)hD_0}{2n}
 \le\frac\theta2 D_0
\]
for all sufficiently large $n$. Consequently, for any sufficiently large
fixed $K$,
\begin{align*}
 &\Exp\left[\exp\left(\frac{L^2+sD_0^2}{2n}\right)
          \boldsymbol1_{\{K\sqrt n<D_0\le h\}}\right]\\
 &\hspace{8mm}\le\ee^{-\theta K\sqrt n/2}\Exp\ee^{\theta D_0}
 =\exp\left(\{C(\ee^\theta-1)-\theta K/2\}\sqrt n+O(1)\right)
 \longrightarrow0.
\end{align*}
On $D_0\le K\sqrt n$, the integrand is bounded and converges in
probability to $\exp((1+s)C^2/2)$. Thus the expectation in
\eqref{eq:secondpoisson} tends to $\exp(3C^2/4)$.
Combining this with \eqref{eq:halfprob} proves \eqref{eq:halfmain}.

\begin{remark}[The two sources of the boundary constant]
The term $3C^2/4$ is the combined curvature correction from the rising
binomial product and the depleted factorial of the large part. The term
$C\sqrt{\nu+2}$ is the contribution of size-two small parts. Omitting
either mechanism gives a wrong multiplicative equivalent. The
normalization shift $\nu$ changes the second contribution, but not the
first.
\end{remark}

\section{Fixed-radius stability and exact coefficient type}
\label{sec:type}

\subsection{A sharp fixed-ball threshold}
\begin{corollary}[Fixed-radius boundedness]\label{cor:fixedball}
For fixed $s,A,C,\nu$, there is a finite $C'$ such that
\[
 f\in\mathcal B_{s,A,C}^{(\nu)}
 \quad\Longrightarrow\quad
 |[z^n]f\inv|\le C' A^{n-1}w_{n-1}\quad(n\ge2)
\]
if and only if $s\ge1$.
\end{corollary}
\begin{proof}
The assertion is equivalent to boundedness of $R_n$. For $s\ge1$,
\eqref{eq:phase} supplies a finite limit, and each of the finitely many
initial $R_n$ is finite. For $0<s<1$, the single extremizer $F_*$ has
$R_n\to\infty$.
\end{proof}

The threshold does not mean that Gevrey inversion fails below order one.
For every $s>0$ and every $A'>A$, \eqref{eq:coarsecost} implies a bound
at radius parameter $A'$ with a constant independent of $n$. What fails
below one is retaining the \emph{same} $A$ with only a constant prefactor.

\subsection{Type preservation, included as a tool}
For $f(z)=\sum f_nz^n$, define
\begin{equation}
 \typ(f)=\limsup_{n\to\infty}
     \left(\frac{|f_n|}{(n!)^s}\right)^{1/n}\in[0,\infty].
 \label{eq:type}
\end{equation}
The convention $0^{1/n}=0$ is understood. Finite initial coefficients do
not affect this value. Replacing $(n!)^s$ by $((n-1)!)^s$ does not alter
it either, since $n^{s/n}\to1$.

\begin{proposition}[Exact inverse type; a known repository result]
\label{prop:inverse-type}
For $s>0$ and $f(z)=az+O(z^2)$ with $a\ne0$,
\begin{equation}
 \typ(f\inv)=\frac{\typ(f)}{|a|}.
 \label{eq:inversetype}
\end{equation}
The equality includes zero and infinite types.
\end{proposition}
\begin{proof}
First suppose $a=1$ and $\typ(f)<\infty$. For every $A>\typ(f)$,
there is a finite $C$ with $|f_{j+1}|\le CA^j(j!)^s$. This follows
from the definition of limsup; the change of index and the polynomial
factor $(j+1)^s$ can be absorbed in the strict exponential margin.
The majorant \eqref{eq:coarsecost} gives $\typ(f\inv)\le A$.
Let $A\downarrow\typ(f)$, including arbitrarily small positive $A$
when the type is zero. Applying the same inequality to $f\inv$ proves
the reverse inequality. If $f$ has infinite type but its inverse had
finite type, this finite-type argument applied to the inverse would be
a contradiction.

For general $a$, let $F=f/a$. It is tangent to the identity and has the
same type as $f$. Since $f\inv(w)=F\inv(w/a)$, input scaling gives
\eqref{eq:inversetype}.
\end{proof}

The general log-convex theorem in \cite{weighted} is broader than
\cref{prop:inverse-type}. The point of the present proof is that it
exhibits the precise information that the type discards.

\subsection{Borel radii, without a summability claim}
For the coefficient normalization
\[
 \mathcal B_s f(\xi)=\sum_{n\ge0}\frac{f_n}{(n!)^s}\xi^n,
\]
Cauchy--Hadamard gives radius $1/\typ(f)$, with the usual conventions at
zero and infinity. The frequently used normalization
\[
 \mathfrak B_s f(\xi)=\sum_{n\ge0}\frac{f_n}{\Gamma(1+sn)}\xi^n
\]
has radius $s^s/\typ(f)$, because Stirling's estimate gives
\[
 \left(\frac{(n!)^s}{\Gamma(1+sn)}\right)^{1/n}\longrightarrow s^{-s}.
\]
Neither radius determines analytic continuation, growth along a ray, or
the existence of a Laplace integral. None of the extremal results above
proves Borel summability or resurgence of the extremizer or of an
exponential-feedback series.

\section{Composition of two divergent series}\label{sec:composition}

\begin{theorem}[A sharp maximum law]\label{thm:composition}
Fix $s>0$. Let $f(z)=az+O(z^2)$ and $g(z)=bz+O(z^2)$, where
$a,b\ne0$. Then
\begin{equation}
 \typ(f\circ g)\le\max\{\typ(g),\ |b|\typ(f)\}.
 \label{eq:compositionbound}
\end{equation}
If $\typ(g)\ne |b|\typ(f)$, equality holds. The equality assertion
also applies when exactly one of the competing types is infinite.
\end{theorem}

\subsection{The finite composition majorant}
We first prove the inequality for tangent-to-identity inputs of finite
type. Choose $A>\typ(f)$ and $B>\typ(g)$ and constants $C,D>0$ with
\begin{equation}
 |f_k|\le CA^{k-1}((k-1)!)^s,\quad
 |g_k|\le DB^{k-1}((k-1)!)^s\qquad(k\ge2).
 \label{eq:compositioninputs}
\end{equation}
In this section $P_0(m,\ell)$ denotes \eqref{eq:convolution} with
$\nu=0$, and $K=K(s,0)$. Write $g=z(1+v)$.
The terms $k=1$ and $k=n$ in $[z^n]f(g)$ contribute $g_n$ and $f_n$.
In every other term, use $\ell\ge1$ nonlinear factors of $v$ in
$(1+v)^k$. Their total extra degree is $n-k$, and their contribution is
bounded by
\begin{equation}
 CA^{k-1}((k-1)!)^s\binom{k}{\ell}
       D^\ell B^{n-k}P_0(n-k,\ell).
 \label{eq:compositionterm}
\end{equation}
Let $T=\max(A,B)$. By \cref{lem:twoconv},
\[
 P_0(n-k,\ell)\le K^{\ell-1}((n-k-\ell+1)!)^s.
\]
Both factorial arguments $k-1$ and $n-k-\ell+1$ are positive, and
sum to $n-\ell$. Hence
\[
 (k-1)!(n-k-\ell+1)!\le(n-\ell-1)!.
\]
Also $A^{k-1}B^{n-k}\le T^{n-1}$ and
$\sum_k\binom{k}{\ell}\le\binom{n+1}{\ell+1}$.
Summing \eqref{eq:compositionterm} proves
\begin{align}
 |[z^n]f(g)|\le{}&|f_n|+|g_n|+
 C T^{n-1}((n-1)!)^s S_n,\label{eq:compositionfinite}\\
 S_n={}&\sum_{\ell=1}^{n-2}D^\ell K^{\ell-1}
 \binom{n+1}{\ell+1}
 \left(\frac{(n-\ell-1)!}{(n-1)!}\right)^s.
 \label{eq:Sn}
\end{align}
This is a finite inequality valid without positivity or convergence.

For $\ell\le n/2$, the factorial ratio is at most
$(n/2)^{-s\ell}$ and the binomial factor at most
$(2n)^{\ell+1}/(\ell+1)!$. For larger $\ell$, the factorial ratio
has at least $\lfloor n/2\rfloor$ factors of order $n^{-s}$, while
all other factors together are at most exponential in $n$. Therefore
\begin{equation}
 S_n\le C_1 n\exp(C_2 n^{\max(1-s,0)})+
               \exp(-c n\log n+O(n)).
 \label{eq:Snbound}
\end{equation}
The right side has $n$th root tending to $1$. Thus
$\typ(f\circ g)\le T$. Sending $A\downarrow\typ(f)$ and
$B\downarrow\typ(g)$ proves the tangent inequality. If an input type
is infinite the upper bound is automatic.

For arbitrary nonzero linear coefficients, put
$\widehat g=g/b$ and $F_b(z)=f(bz)/(ab)$. These are tangent to the
identity, and $f\circ g=ab(F_b\circ\widehat g)$. Since
$\typ(F_b)=|b|\typ(f)$ and $\typ(\widehat g)=\typ(g)$,
\eqref{eq:compositionbound} follows.

\subsection{Why unequal types cannot cancel}
Let $h=f\circ g$, and abbreviate
\[
 X=|b|\typ(f),\qquad Y=\typ(g),\qquad Z=\typ(h).
\]
Apply the already proved inequality to
$f=h\circ g\inv$ and $g=f\inv\circ h$, using
\cref{prop:inverse-type}. One obtains
\begin{equation}
 X\le\max(Y,Z),\qquad Y\le\max(X,Z),\qquad Z\le\max(X,Y).
 \label{eq:threemax}
\end{equation}
If $X\ne Y$, these three inequalities force $Z=\max(X,Y)$.
This also proves the extended-real assertion. It completes the proof
of \cref{thm:composition}.

\begin{example}[Both equality and strict inequality are necessary]
Let
$f(z)=z+\sum_{n\ge2}A^{n-1}((n-1)!)^s z^n$ and define $g$ similarly
with $B\ne A$. Then their types are $A,B$, and the composition has type
$\max(A,B)$. On the other hand, taking $g=f\inv$ gives two equally
typed divergent inputs whose composition is the identity. Its type is
zero. Thus an unconditional equality in \eqref{eq:compositionbound}
would be false.
\end{example}

\begin{corollary}[A type filtration]
For each finite $T\ge0$, the tangent-to-identity formal series with
$\typ(f)\le T$ form a group under composition. If the inner and outer
types are different, the larger type survives composition.
\end{corollary}
\begin{proof}
Use \cref{thm:composition,prop:inverse-type}; the identity has type zero.
\end{proof}

\begin{remark}[Analytic coordinate changes]
For an analytic invertible $\phi(z)=cz+O(z^2)$,
\[
 \typ(\phi\inv\circ f\circ\phi)=|c|\typ(f).
\]
At positive type this follows from the unequal-type equality, since an
analytic series has zero $s$-type for $s>0$. At zero type it follows
from the upper bound. This recovers the factorial specialization of the
coordinate law already present in \cite{weighted}; it is a consequence,
not an additional novelty claim.
\end{remark}

\section{A positive-forward obstruction to inverse root limits}
\label{sec:counterexample}

\begin{theorem}[Full forward roots do not force full inverse roots]
\label{thm:rootcounter}
For every $s>0$ and $A>0$, there is a tangent-to-identity formal series
$F$ with strictly positive coefficients in every positive degree and
\begin{equation}
 \lim_{n\to\infty}
 \left(\frac{[z^n]F}{((n-1)!)^s}\right)^{1/n}=A,
 \label{eq:fullforward}
\end{equation}
whose inverse has infinitely many zero coefficients and satisfies
\begin{equation}
 \liminf_{n\to\infty}
 \left(\frac{|[z^n]F\inv|}{((n-1)!)^s}\right)^{1/n}=0,
 \qquad
 \limsup_{n\to\infty}
 \left(\frac{|[z^n]F\inv|}{((n-1)!)^s}\right)^{1/n}=A.
 \label{eq:sparseroots}
\end{equation}
\end{theorem}
\begin{proof}
Define the sparse series
\begin{equation}
 H(z)=z-\sum_{\substack{n\ge2\\n\text{ even}}}
                  A^{n-1}((n-1)!)^s z^n,
 \qquad F=H\inv.
 \label{eq:sparseH}
\end{equation}
The inverse equation is
\[
 F=z+\sum_{\substack{m\ge2\\m\text{ even}}}
          A^{m-1}((m-1)!)^s F^m.
\]
Its recursive construction has nonnegative coefficients, and the term
$AF^2$ alone forces strict positivity in every positive degree. In
particular $F_2=A$.

For even $n$, the term with $m=n$ yields
\[
 F_n\ge A^{n-1}((n-1)!)^s.
\]
For odd $n\ge3$, use $m=n-1$ and choose exactly one quadratic term in
$F^{n-1}$:
\[
 F_n\ge A^{n-2}((n-2)!)^s(n-1)F_2
       =A^{n-1}(n-1)((n-2)!)^s.
\]
After normalization by $((n-1)!)^s$, the latter lower bound is
$A^{n-1}(n-1)^{1-s}$. Thus both parity subsequences have normalized
root liminf at least $A$.

The sparse series $H$ has type $A$, and
\cref{prop:inverse-type} gives $\typ(F)=A$. The upper limsup for $F$
is therefore $A$, proving \eqref{eq:fullforward}. Since $F\inv=H$,
the odd coefficients vanish and the even normalized roots tend to $A$.
This proves \eqref{eq:sparseroots}.
\end{proof}

This example is deliberately not an exponential-feedback kernel. It
shows exactly what a theorem based only on positivity and a forward
root limit cannot establish. Eventual nonvanishing, a sign law, or a
full root limit for a particular feedback inverse needs additional
structure specific to that kernel.

The restriction $s>0$ throughout the type theory is essential. At
$s=0$, $f=z-z^2$ has type zero, whereas its inverse has coefficients
$\frac1n\binom{2n-2}{n-1}$ and type $4$. Positive factorial growth is
what makes the nonlinear overhead subexponential.

\section{Reproducible calculations and what they establish}
\label{sec:computations}

\subsection{An exact coefficient algorithm}
Let $f(z)=z(1+\sum_{j\ge1}v_jz^j)$ and, for a fixed $n$, write
$(1+\sum v_jz^j)^{-n}=\sum b_mz^m$. Differentiating formally gives
\begin{equation}
 b_0=1,\qquad
 b_m=-\frac1m\sum_{j=1}^m\{m+(n-1)j\}v_jb_{m-j},\qquad
 [z^n]f\inv=\frac{b_{n-1}}n.
 \label{eq:algorithm}
\end{equation}
For integer input, the implementation checks exact divisibility in both
steps rather than silently rounding. Composition is implemented by an
independent truncated polynomial multiplication routine.

The recorded standard-library run used degree $45$ and performed
\textbf{5,338 scalar equalities or inequalities}. It checked both inverse
identities for four factorial models, all $512$ sign assignments in a
degree-ten coefficient box, the sparse-inverse counterexample and its
lower bounds, and the identity
$(f\circ g)\inv=g\inv\circ f\inv$ for two factorially divergent
inputs. These are finite exact checks; they do not certify a limit.

For $s=C=A=1$ and $\nu=0$, the initial coefficients of $G_*$ are
\[
 1,\ 1,\ 4,\ 21,\ 128,\ 862,\ 6264,\ 48549,\ 399352,\ 3484746.
\]
For the sparse inverse in \eqref{eq:sparseH} at $s=A=1$, the forward
coefficients begin
\[
 1,\ 1,\ 2,\ 11,\ 50,\ 330,\ 1956,\ 15459,\ 109322,\ 1040414.
\]
The complete recorded coefficients and checks are in
\texttt{data/recorded/verification.json}.

\subsection{Positive-term diagnostics, not interval certificates}
For the extremizer with $A=1$, normalize the recurrence by $w_m$.
Writing $B_m=b_m/w_m$ for $m\ge1$ gives
\begin{equation}
 B_m=Cn+\frac Cm\sum_{j=1}^{m-1}\{m+(n-1)j\}
              \frac{w_jw_{m-j}}{w_m}B_{m-j},\qquad
 R_n=\frac{B_{n-1}}{nC}.
 \label{eq:scaledalgorithm}
\end{equation}
This avoids building the final huge factorial before normalization.
All summands are positive. The diagnostic implementation uses
extended-range \texttt{longdouble}, rejects environments where that
format has only ordinary double range, and checks finiteness. Three
cases at degree $80$ were independently recomputed at $80$ decimal digits
with \texttt{mpmath}; relative differences were below $5\cdot10^{-18}$.
Neither computation uses directed rounding.

\begin{table}[htbp]
\centering\small
\begin{tabular}{rrrrr}
\toprule
$s$ & $n$ & $\log R_n$ & $\log R_n/n^{1-s}$ & $\log R_n-n^{1-s}$\\
\midrule
0.25 & 320 & 193.909447 & 2.562928 & 118.250119 \\
0.25 & 1280 & 363.057932 & 1.696557 & 149.061035 \\
0.5 & 320 & 20.613024 & 1.152303 & 2.724481 \\
0.5 & 1280 & 38.185907 & 1.067329 & 2.408819 \\
0.75 & 320 & 4.404885 & 1.041471 & 0.175399 \\
0.75 & 1280 & 6.062521 & 1.013563 & 0.081126 \\
1 & 320 & 1.015916 & 1.015916 & 0.015916 \\
1 & 1280 & 1.003924 & 1.003924 & 0.003924 \\
1.25 & 320 & 0.238729 & 1.009699 & 0.002293 \\
1.25 & 1280 & 0.167556 & 1.002219 & 0.000371 \\
\bottomrule

\end{tabular}
\caption{Recorded diagnostics for $C=1$, $\nu=0$. Every displayed entry is a rounded floating-point diagnostic. At $s=1/2$, the last column tends to $3/4+\sqrt2=2.16421356\ldots$; convergence is visibly not complete at the displayed degrees.}
\label{tab:diagnostics}
\end{table}

At $s=1/4$, the ratio in the fourth column is still about $1.70$ at
$n=1280$, although its proved limit is $1$. This is an important caution:
a leading logarithmic theorem need not provide a useful small-degree
approximation. We do not use the table to assert an onset or error rate.

\begin{table}[htbp]
\centering\small
\begin{tabular}{rrrr}
\toprule
$t$ & $s=1+t/\log1280$ & observed $R_{1280}$ & $\exp(\ee^{-t})$\\
\midrule
-2 & 0.720460 & 1823.607153 & 1618.177992 \\
-1 & 0.860230 & 15.460266 & 15.154262 \\
0 & 1.000000 & 2.728969 & 2.718282 \\
1 & 1.139770 & 1.446060 & 1.444668 \\
2 & 1.279540 & 1.145256 & 1.144921 \\
\bottomrule

\end{tabular}
\caption{The order-one crossover with $C=1$, $\nu=0$. The rightmost column is the limit in \eqref{eq:crossover}, not a finite-degree enclosure.}
\end{table}

\subsection{Rebuilding without modifying recorded results}
The package has no network dependency. From its root directory:
\begin{quote}\small\ttfamily
python code/verify.py --order 45\\
python code/verify.py --order 45 --diagnostics --mp-check\\
pdflatex -interaction=nonstopmode -halt-on-error article.tex\\
pdflatex -interaction=nonstopmode -halt-on-error article.tex
\end{quote}
The exact tests need only Python's standard library. Diagnostics require
NumPy, and the high-precision cross-check requires mpmath. The default
output directory is \texttt{data/rerun/}, so an ordinary rerun does not
overwrite the recorded evidence. The numerical tables used by this
article are supplied as small TeX files generated from the recorded CSVs.

\section{Implications for the transseries program}\label{sec:implications}

\subsection{What the result adds below type preservation}
A type identifies a coefficient series only up to factors whose $n$th
roots tend to $1$. At fixed positive Gevrey order, that equivalence is
too coarse to control a critical-radius coefficient norm. The extremal
cost quantifies the missing information. In the shifted normalization,
quadratic coefficients generate approximately $Cn^{1-s}$ small
Lagrange blocks, and their accumulation produces $\exp(Cn^{1-s})$.
At and above $s=1$ their number stays bounded or tends to zero.

This has a practical implication for certified symbolic manipulation:
retaining an optimal exponential constant does not by itself justify
retaining a fixed coefficient prefactor. For $s<1$, a program that
returns the same $A$ and a constant-size majorant after reversion will
fail on the explicit extremizer. It must either track a subexponential
envelope, allow an arbitrarily small increase of $A$, or exploit extra
structure of its particular input.

\subsection{What transfers to countable exponential feedback}
The source weighted-type theorem is motivated by
\[
 U(q)=\sum_{j\ge1}c_jq^j\exp(\lambda_jU(q)).
\]
Once a separately established theorem places such a forward series in
a particular factorial coefficient ball, the exact finite majorant
\eqref{eq:exactR} bounds its inverse. The bound is universal and requires
no information about the inverse signs. The composition theorem also
controls the type when two divergent formal correction maps are nested.

These statements do \emph{not} say that a particular feedback inverse
attains the extremal amplification. The extremizer has negative
nonlinear input coefficients; positive-feedback inputs can exhibit
substantial cancellation after inversion. Consequently, this article
makes no new signed equivalent, least-term error law, or summability
claim for the feedback family. Its contribution is the universal
fixed-ball estimate and the sharp filtered composition law.

\subsection{Formal and analytic layers remain separate}
Everything in the proofs before the numerical section concerns ordinary
formal coefficient degrees. An exponential transseries with accumulating
actions may not have a finite degree decomposition of the kind used in
\eqref{eq:convolution}. Applying these arguments to such supports requires
an additional finite-coefficient or summability theorem. Similarly,
a formal inverse need not be the asymptotic expansion of a chosen
analytic inverse without an analytic realization theorem. Neither issue
is removed by the positivity of the extremal majorant.

\section{Further research questions and concrete next targets}
\label{sec:questions}

\subsection{Complete multiplicative cost below one half}
For $0<s<1/2$, determine a finite explicit expansion of $\log R_n$
whose error is $o(1)$. The Poisson representation suggests a hierarchy
of scales $n^{1-rs}$. Small parts of size $r$ become visible at
$s=1/r$, but product curvature and tilted fluctuations also contribute.
The formula for $s=1/2$ proves that simply summing the small-part
intensities is insufficient. A useful next theorem would give the
complete equivalent for $1/3<s<1/2$, with its constant at $s=1/3$.
No coefficient for that further expansion is asserted here.

\subsection{Uniform transitions at the lower thresholds}
Construct analogues of \eqref{eq:crossover} around $s=1/r$ for
$r\ge2$. The variables $(s-1/r)\log n$ are natural candidates, but
one must first subtract all larger contributions. A successful theorem
should specify both a uniform remainder and the normalization under
which a finite limiting profile appears.

\subsection{Replacing factorials by a local-quotient theory}
For a log-convex weight $M$, seek hypotheses under which the critical
inverse cost is controlled sharply by $nM_{n-1}/M_n$, after the
appropriate index convention is fixed. In the present model the ratio
is asymptotic to $n^{-s}$, explaining the scale $n^{1-s}$.
Slowly varying perturbations, abrupt quotient changes, and factorial
weights with logarithmic deficits need separate uniformity analysis.
The repository's general zero-loss theorem alone is not enough to
identify their extremal subexponential constants.

\subsection{Two divergent inputs for general weights}
Extend \cref{thm:composition} beyond factorial powers. A candidate
hypothesis is a weighted two-part convolution estimate with
subexponential constants uniform in the number of parts. Determine
whether log-convexity and superexponential root growth alone suffice,
or construct a counterexample. Equal-type cancellation must remain
part of the statement.
% ed. (2026-09-29): answered by the companion package
\begin{ednote}
Answered by \path{Exact_Weighted_Type_Beyond_Log_Convexity/}: log-convexity
and superexponential root growth suffice, and more generally so does
subexponential equivalence to a supermultiplicative weight with
superexponential roots, with equal-type cancellation filling all of
$[0,T]$; see the note at the end of Section~\ref{sec:target}.
\end{ednote}

\subsection{Typical inputs versus the extremizer}
Classify coefficient sign patterns or asymptotic regularity assumptions
under which $|[z^n]f\inv|$ is asymptotic to a fixed fraction of the
extremal coefficient. An opposite target is a sharp suppression theorem
for positive forward coefficients. The sparse example shows that a
root-limit hypothesis by itself cannot provide such a classification.

\subsection{Exact remainder bounds at finite degree}
Turn the localization proof into a computable error certificate for
\eqref{eq:multmain} and \eqref{eq:halfmain}, with an explicit optimal
cutoff $h$ and a rigorous bound for every remainder. The finite sum
\eqref{eq:exactR} already gives an exact extremum, but the present
asymptotic proof does not supply an optimized, practical onset bound.
Directed-rounding evaluation of the positive recurrence would provide
a complementary finite-degree certificate.

\subsection{Several variables and operator-valued coefficients}
Define a coefficient ball for homogeneous polynomial maps or bounded
multilinear operators with an invertible linear part. Determine which
normed extremal problem replaces the one-variable negative series and
whether the threshold $s=1$ persists without dimension-dependent
exponential loss. A type-level theorem and a sharp fixed-ball theorem
should be treated as different objectives.
% ed. (2026-09-29): the type-level part is in the companion package
\begin{ednote}
The type-level statement for tangent maps on Banach spaces, with the
multilinear operator norm and no dimension loss, and sharp two-sided bounds
under a general invertible derivative, are proved in
\path{Exact_Weighted_Type_Beyond_Log_Convexity/} (its
Theorems~\texttt{thm:banach} and \texttt{thm:jacobian}). The fixed-ball
problem posed here remains open.
\end{ednote}

\subsection{A realistic formalization sequence}
A Lean formalization can begin with the finite identities
\eqref{eq:lagrange}, \eqref{eq:exactR}, and \eqref{eq:compositionfinite}.
The next independent component is the factorial concentration and
convolution library, including a summable uniform bound for $K$.
A third component proves the tail estimates and the type maximum law.
Only after those finite and limsup statements are established is it
reasonable to formalize localization, the Poisson normalization, and
uniform asymptotics. The present Python checks should remain explicitly
separate from that program.

\section{Conclusion}
The Gevrey version of the repository's optimal-overhead question has a
sharp answer at the leading subexponential scale: a single extremizer
attains $\log R_n\sim Cn^{1-s}$ for $0<s<1$. Its full multiplicative
behavior is determined here at and above one half, including the
boundary constant, while the uniform transition at order one separates
fixed-radius boundedness from type preservation. Two divergent formal
maps obey a maximum law at the level of Gevrey type, and unequal types
cannot cancel. Full root limits require additional information, as the
explicit positive-forward counterexample demonstrates.

These are extensions of existing repository results, not replacements
for the already established zero-loss theorem. They isolate what must
be tracked beyond an optimal exponential constant and supply finite
formulas suitable for both implementation and subsequent formalization.

\appendix
\section{A compact proof audit}\label{app:audit}
The logical dependencies can be checked without reading the numerical
section. Formal Lagrange inversion implies the exact extremizer formula.
The two-part convolution lemma implies the length tail bound. Capped
factorial concentration plus that length tail implies localization.
Localization and an elementary negative-binomial upper bound, paired
with the all-size-one lower bound, prove the sharp logarithmic cost.
Localization at a shrinking cutoff and a finite Poisson normalization
prove the multiplicative equivalent. The order-one-half correction uses
one extra product expansion and a separately proved exponential
integrability estimate. No empirical observation is used in these steps.

The type-invariance proposition uses only the coarse convolution bound
and inversion in both directions. The composition theorem uses the
finite composition majorant for its upper bound and type invariance
for its noncancellation conclusion. The root-limit counterexample uses
positivity, two explicit coefficient lower bounds, and type invariance.

The potentially delicate points are therefore exposed: uniformity in
the convolution constant; entropy control with $t\log(\ee n/t)$ rather
than the too-coarse $t\log n$ in the fixed-fraction localization proof;
uniqueness of the large part; the distinction between an absolute
$o(1)$ discarded contribution and a relative error; and the
exponential-integrability argument at $s=1/2$. These points are all
proved explicitly in the main text.

\section{Package contents and provenance}\label{app:package}
\begin{longtable}{>{\raggedright\arraybackslash}p{.39\textwidth}p{.54\textwidth}}
\toprule
File & Purpose\\
\midrule
\endhead
\texttt{article.tex}, \texttt{article.pdf} & Source and compiled article.\\
\texttt{README.md} & Scope, main results, and rebuild commands.\\
\texttt{code/verify.py} & Exact finite checks and optional diagnostics.\\
\texttt{code/render\_tables.py} & Regenerate the printed tables from CSV files.\\
\path{data/recorded/verification.json} & Recorded exact-check ledger and high-precision comparisons.\\
\texttt{data/recorded/diagnostics.csv} & Positive-term extremal computations.\\
\texttt{data/recorded/crossover.csv} & Order-one transition diagnostics.\\
\texttt{data/*\_table.tex} & Tables generated from the recorded CSVs.\\
\texttt{data/build\_validation.json} & Build and layout validation record.\\
\texttt{provenance.json} & Repository snapshot, source paths, and dependency boundaries.\\
\bottomrule
\end{longtable}
No repository source file is changed by this package. No claim is made
that the full ProveIt corpus was exhaustively searched. In particular,
concurrently arriving or unintegrated archives are not certified to be
disjoint from this article. The explicit crosswalk in \cref{sec:target}
records the actual target and avoids claiming the already settled
inverse-type question as a new result.

\begin{thebibliography}{99}
\bibitem{repo}
V. Reshetnikov, \emph{ProveIt}, repository snapshot, inspected
29 September 2026.\par
\commit.\par
\href{https://github.com/VladimirReshetnikov/ProveIt/commit/5c695cfdf70a8b6c91bd5b2c1e99e3baecd5a864}{GitHub snapshot}.

\bibitem{inventory}
ProveIt contributors, \emph{Series and transseries}, directory README
at the snapshot in \cite{repo}, including the arrival inventory and
editorial scope notes.
\href{https://github.com/VladimirReshetnikov/ProveIt/blob/5c695cfdf70a8b6c91bd5b2c1e99e3baecd5a864/Analysis/Transseries/docs/series-and-transseries/README.md}{Source inventory}.

\bibitem{weighted}
ProveIt research package, \emph{Sharp Weighted Type Under Formal
Reversion: zero-loss inversion, exact Borel radii, and non-D-finiteness
in countable exponential feedback}, 2026, conventional research
manuscript with editorial notes; especially the further-research
subsections on optimal subexponential cost and two divergent inputs.
\href{https://github.com/VladimirReshetnikov/ProveIt/blob/5c695cfdf70a8b6c91bd5b2c1e99e3baecd5a864/Analysis/Transseries/docs/series-and-transseries/Sharp_Weighted_Type_Formal_Reversion/article.tex}{Pinned TeX source}.

\bibitem{amplitude}
ProveIt research package, \emph{Amplitude--Slope Compensation in
Feedback Transseries: sharp regularity, negative-ray summation, and
damping-induced condensation}, 2026. Its inverse-type question is
cross-referenced as already answered by \cite{weighted} in
\cite{inventory}. Used here only to identify this overlap, not as a
mathematical input to any proof.
\href{https://github.com/VladimirReshetnikov/ProveIt/blob/5c695cfdf70a8b6c91bd5b2c1e99e3baecd5a864/Analysis/Transseries/docs/series-and-transseries/Amplitude_Slope_Compensation_Feedback_Transseries/article.tex}{Pinned source}.

\bibitem{gessel}
I. M. Gessel, Lagrange inversion, \emph{Journal of Combinatorial Theory,
Series A} \textbf{144} (2016), 212--249.
\href{https://doi.org/10.1016/j.jcta.2016.06.018}{doi:10.1016/j.jcta.2016.06.018};
\href{https://arxiv.org/abs/1609.05988}{arXiv:1609.05988}.

\bibitem{borinsky}
M. Borinsky, Generating asymptotics for factorially divergent sequences,
\emph{Electronic Journal of Combinatorics} \textbf{25}(4) (2018), P4.1.
\href{https://doi.org/10.37236/5999}{doi:10.37236/5999};
\href{https://arxiv.org/abs/1603.01236}{arXiv:1603.01236}.

\bibitem{rscomposition}
A. Rainer and G. Schindl, Composition in ultradifferentiable classes,
\emph{Studia Mathematica} \textbf{224}(2) (2014), 97--131.
\href{https://doi.org/10.4064/sm224-2-1}{doi:10.4064/sm224-2-1};
\href{https://arxiv.org/abs/1210.5102}{arXiv:1210.5102}.

\bibitem{rsstability}
A. Rainer and G. Schindl, Equivalence of stability properties for
ultradifferentiable function classes, \emph{Revista de la Real Academia
de Ciencias Exactas, Fisicas y Naturales. Serie A. Matematicas}
\textbf{110} (2016), 17--32.
\href{https://doi.org/10.1007/s13398-014-0216-0}{doi:10.1007/s13398-014-0216-0};
\href{https://arxiv.org/abs/1407.6673}{arXiv:1407.6673}.
\end{thebibliography}
\end{document}
